% Generated draft: requires mathematical and source review.
\documentclass[11pt]{article}
\usepackage[margin=25mm]{geometry}
\usepackage{fontspec}
\setmainfont{Latin Modern Roman}
\usepackage{amsmath,amssymb,mathtools,unicode-math}
\setmathfont{Latin Modern Math}
\usepackage{enumitem,needspace,xurl,hyperref}
\setlength{\parindent}{0pt}
\setlength{\parskip}{5pt}
\setlength{\emergencystretch}{3em}
\newcommand{\N}{\mathbb{N}}
\newcommand{\Z}{\mathbb{Z}}
\newcommand{\Q}{\mathbb{Q}}
\newcommand{\R}{\mathbb{R}}
\newcommand{\C}{\mathbb{C}}
\newcommand{\F}{\mathbb{F}}
\newcommand{\A}{\mathbb{A}}
\newcommand{\PP}{\mathbb{P}}
\newcommand{\E}{\mathbb{E}}
\newcommand{\eps}{\varepsilon}
\newcommand{\dd}{\,\mathrm{d}}
\begin{document}
\section*{UnsolvedMath ID: 1900022}
\textbf{Source problem code:} catalogue entry
\subsection{Definitions and mathematical statement}
A knot is a tame embedding of $S^1$ in $S^3$, considered up to ambient isotopy. Its bridge number $b(K)$ is the minimum number of local maxima of a generic height function over representatives in $\R^3\subset S^3$. A knot of bridge number exactly three is a 3-bridge knot; some terminology also allows knots admitting a three-bridge presentation, which means $b(K)\leq3$. The supplied source does not specify which convention is intended.

A finite ordinary continued fraction has the form
\[
 [a_0;a_1,\ldots,a_r]
 =a_0+\frac1{a_1+\dfrac1{\ddots+\dfrac1{a_r}}},
\]
with integer entries and all expressions defined. The problem asks for a generalization of continued fractions that describes 3-bridge knots, extending the continued-fraction treatment of 2-bridge knots.

The supplied assessment interprets the goal as a parametrization reaching all 3-bridge knots. However, the record does not prescribe the generalized fraction's syntax, geometric construction, allowed equivalence moves, or whether ``describe'' requires only surjective representation, recognition of isotopic outputs, a canonical classification, or formulas for knot invariants. These are distinct possible goals. Thus no particular algorithm or unproved completeness criterion is silently substituted for the requested generalization.
\subsection{Short English statement}
Develop a continued-fraction framework for 3-bridge knots. The supplied assessment labels this partially solved: it describes the 2-bridge case as understood through continued-fraction parametrizations and Jones-polynomial formulas, and mentions Chebyshev/billiard-diagram and cluster-algebra approaches for 3-bridge knots. It reports no found continued-fraction generalization parametrizing all such knots and calls that part open. No specific supporting bibliography or precise completeness theorem is supplied, so these are attributed review observations rather than an independently established current-status determination.
\subsection{Sources}
\begin{itemize}
\item UnsolvedMath permanent ID 1900022, source code catalogue entry, ``Geometry of Continued Fractions --- Further open questions.'' Registry curation and metadata: Ulam AI and UnsolvedMath Contributors, CC BY 4.0; updated 2026-08-06.
\item Karpenkov, \emph{Open problems in geometry of continued fractions} (2017), Problem 22: \url{https://arxiv.org/abs/1712.01450}. Supplied source-TeX extraction accessed 2026-07-29. The metadata says the source presents an open item and marks status/release and rights review NEEDS\_REVIEW.
\item The supplied progress summary names continued-fraction Jones formulas, Chebyshev/billiard diagrams and cluster-algebra methods without identifying specific papers or theorem statements. No missing reference is invented.
\end{itemize}
\end{document}
